\section{Research Gap and Contribution}
\label{sec:research_gap}

Previous research links institutional ownership, crowded positions, and liquidity-driven trading to non-fundamental price pressure, stock-price fragility, and downside tail risk \parencite{coval2007asset,greenwood2011stock,brown2022crowded}. More recent work has begun to model institutional holdings through dynamic graphs. \Textcite{izadifar2026institutional} represent Form 13F holdings as a temporal bipartite network of institutional managers and securities and use graph-learning methods to forecast future portfolio allocations. \Textcite{tagndynamic2026} instead incorporate institutional co-holdings into a broader multi-relational graph, alongside price correlations and supply-chain relationships, to model financial-risk transmission and predict extreme volatility. However, these studies do not establish whether institutional ownership topology, considered as a distinct source of information, improves the prediction of stock-level downside risk beyond market characteristics and conventional crowding measures. \Textcite{izadifar2026institutional} predict future holdings rather than risk outcomes, while \textcite{tagndynamic2026} predict two-sided extreme volatility using several network layers simultaneously. The contribution of institutional co-holdings therefore cannot be isolated from the model's other information channels. This leaves open the incremental predictive value of ownership topology for downside risk and the economic relevance of the resulting predictions.

This thesis addresses the gap by representing institutional holdings as a dynamic manager--stock ownership network and testing whether its topology improves future downside-risk predictions beyond market characteristics and conventional ownership measures. It also examines whether the resulting predictions contain an out-of-sample investment signal. In doing so, the thesis makes three related contributions: empirically, it tests whether conventional ownership measures and ownership-network characteristics provide incremental predictive information; methodologically, it compares tabular models using market, conventional ownership, and engineered network features with graph neural networks that learn directly from the ownership structure; and economically, it transforms the selected model's predictions into a stock-level Fragility Score and evaluates its relevance through downside-risk sorting and out-of-sample portfolio tests.